\documentclass[../main]{subfiles}
\begin{document}
\chapter{Ecuaciones de continuidad}
\img{images/04/bernoulli.jpg}{Caudal continuidad}
\info{
Contenido}{
  Ecuaciones de continuidad.
}

\info{Caudal}{
  Se define caudal o gasto de una corriente líquida al volumen que
atraviesa una sección de la misma en la unidad de tiempo.}

Sea un trozo de filamento de corriente con sección transversal $S$,
por la cual pasan partículas líquidas perpendicularmente con una
velocidad $v$. En un intervalo de tiempo $\Delta t$ un volumen líquido.
\begin{equation}
  \Delta V = S \cdot v \Delta t
\end{equation}

Por definición, el caudal elemental que circula por la sección del
filamento de corriente será:

\begin{equation}
  Q= \dfrac{\Delta V}{\Delta t}
\end{equation}

En consecuencia el caudal total que pasa a través de la sección $S$
(plana y perpendicular a la dirección del movimiento, de una
  corriente líquida formada por infinitos filamentos de corriente valdrá.
  \begin{equation}
    Q= v \cdot S
  \end{equation}

  \section{Ecuación de continuidad}
  Si no hay pérdidas de fluido dentro de un tubo
  uniforme, la masa de fluido que entra es igual a
  la que sale en un tubo en un tiempo. Por
  ejemplo, en la figura, la masa ($\Delta m_1$) que entra y la masa
  ($\Delta m_2$) que salen en el tubo durante
  un tiempo corto ($\Delta t$) son

  $\Delta m_1 = \rho_1 \Delta V_1 = \rho_1A_1\Delta x_1 = \rho_1 A_1
  v_1 \Delta t$

  $\Delta m_2 = \rho_1 \Delta V_2 = \rho_1 A_2 \Delta x_2 = \rho_1
  A_2 v_2 \Delta t$

  Dónde $A_1$ y $A_2$ son las áreas transversales del tubo en la
  entrada y salida. Puesto que la masa se conserva ($\Delta m_1 =
  \Delta m_2$) se tiene

  $\rho_1 A_1 v_1\Delta t = \rho_1 A_2 v_2\Delta t,$

  Esté resultado se se denomina ecuación de continuidad.
  \section{Fluidos ideales}
  Si un fluido es incompresible, su densidad $\rho$ es
  constante, así que

  $A_1 v_1 = A_2 v_2,$

  Dónde dicho resultado se conoce como la ecuación de continuidad un
  fluido ideal.
  Por otro lado la cantidad $A v$ se conoce como caudal promedio $Q$
  y representa el volumen del fluido que pasa por un punto en el tubo
  por unidad de tiempo.
  Nota: la definición formal para el caudal es $Q =\dfrac{\Delta V}{\Delta t}$

  \actividad{ Resolver}
  \begin{enumerate}
    \item Un caudal constante de 20 litros por segundo se derrama
      desde un tanque cilíndrico de 1m de diámetro, a través de un
      conducto circular de 0,1m de diámetro que termina en una
      boquilla de 5cm de diámetro. Determinar la velocidad de la
      corriente líquida en las secciones $d_1; d_2$ y a la salida de
      la boquilla.

      %\includestandalone{fig/tikz/tp4p1}

    \item Un tanque tricapa de $V=650L$ litros lleno de aceite, se
      vacía con caudal constante en $t=12h$ horas. Calcular el tamaño
      de la abertura, si sale con velocidad $v_2=0,01m/s$.
    \item Una cañería de $d_1=1m$ se reduce a $d_2=0.5m$,si la
      velocidad inicial es de $v_1=10m/s$. Calcular la velocidad final.
    \item Un balde $V=5L$ se llena con agua de la canilla de $1/2"$
      en $t=1 min$. Calcular la velocidad $v_2$ que alcanzaría el
      agua si se reduce la mitad su sección.
    \item El caudal de un río es $Q=11\frac{m^3}{s}$ si tiene
      $l=10m$ ancho y $h=2m$ de profundidad. Calcular la velocidad de
      circulación del agua, y además determinar cuanto aumentará su
      velocidad si se disminuye la sección del río en un 30\%.
    \item La cañería forzada de una central hidroeléctrica tiene una
      sección circular de diámetro de $d_1=2m$, si tiene una
      velocidad de $v_1=2\frac{m}{s}$ y se reduce al $d_2=0,5m$ al
      llegar a los álabes de una turbina, calcular la velocidad con
      que impacta el chorro de agua en la turbina.

    \item Un tanque cilíndrico de 2 metros de diámetro se vacía a
      través de un tubo de 0.5 metros de diámetro en 10 segundos.
      Calcula la velocidad de salida del líquido.

    \item Un recipiente cónico se llena con agua a una velocidad
      constante de 0.2 metros cúbicos por segundo. Si el diámetro de
      la base del cono es de 4 metros y la altura del cono es de 6
      metros, calcula la velocidad del agua en la parte más estrecha del cono.

    \item Un conducto con forma de V tiene una sección transversal de
      2 metros cuadrados en la parte superior y una sección
      transversal de 0.5 metros cuadrados en la parte inferior. Si se
      bombea agua a una velocidad de 1 metro cúbico por segundo en la
      parte superior del conducto. Calcular las velocidades inicial y final.

    \item Un tanque cilíndrico de 1 metro de diámetro se llena con
      agua a través de un tubo de 0.2 metros de diámetro en 5
      minutos. Calcula la velocidad de entrada del agua.

    \item Un canal rectangular tiene un ancho de 3 metros y una
      altura de 2 metros. Si se envía agua a un caudal de 4 metros
      cúbicos por segundo, calcula la velocidad del agua en el canal.

    \item Un tanque cilíndrico de 1 metro de diámetro se llena con
      agua a través de un tubo de 0.5 metros de diámetro en 10
      segundos. Calcula la velocidad de entrada del agua.

    \item Un conducto con forma de V tiene una sección transversal de
      2 metros cuadrados en la parte superior y una sección
      transversal de 1 metro cuadrado en la parte inferior. Si se
      bombea agua a una velocidad de 0.5 metros cúbicos por segundo
      en la parte superior del conducto, calcula la velocidad del
      agua en la parte inferior.

    \item Un canal rectangular tiene un ancho de 4 metros y una
      altura de 3 metros. Si se envía agua a un caudal de 6 metros
      cúbicos por segundo, calcula la velocidad del agua en el canal.

    \item Un sistema de tuberías consta de tres tubos conectados en
      serie. El primer tubo tiene un diámetro de 10 cm y una longitud
      de 5 m, el segundo tubo tiene un diámetro de 8 cm y una
      longitud de 3 m, y el tercer tubo tiene un diámetro de 6 cm y
      una longitud de 4 m. El caudal en la entrada del sistema es de
      0.5 m³/s. Calcula la velocidad del flujo en cada parte.

    \item Un tanque de almacenamiento en forma de prisma rectangular
      tiene una base de 4 m de ancho, 6 m de largo y una altura de 3
      m. El tanque se llena con agua a una velocidad constante de 0.1
      m³/s a través de una tubería de 0.2 m de diámetro. Calcula el
      tiempo requerido para llenar el tanque por completo.

    \item Una cuba para la producción de refrigerante para torno de
      2000litros se llena con agua a través de una cañería de 1/2
      pulgada de diámetro, a una velocidad de 1m/s, simultáneamente
      se llena con aceite de corte soluble ACS a través de una
      cañería de 3/4 de pulgada a una velocidad de 0.5m/s, calcular
      el tiempo de llenado del tanque y cuántos litros de agua y de
      aceite ACS completará la mezcla.

  \end{enumerate}

  \end{document}
